% !TEX program = xelatex
\documentclass[10pt,a4paper]{article}

\usepackage{geometry}
\geometry{top=1.2cm, bottom=1.2cm, left=1.4cm, right=1.4cm}

\usepackage{xcolor}
\usepackage{tikz}
\usepackage{enumitem}
\usepackage{tabularx}
\usepackage{array}
\usepackage{fontawesome5}
\usepackage[most]{tcolorbox}
\usepackage{graphicx}
\usepackage{eso-pic}

\definecolor{primary}{HTML}{262E40}
\definecolor{muted}{HTML}{565656}
\definecolor{lightgray}{HTML}{D3D3D3}
\definecolor{headlight}{HTML}{CED4D9}
\definecolor{linkblue}{HTML}{1E6FBA}

\usepackage[colorlinks=true, linkcolor=linkblue, urlcolor=linkblue]{hyperref}

\usepackage{xepersian}
\settextfont[
  Path=/home/ali/.local/share/fonts/,
  Extension=.ttf,
  UprightFont=Vazirmatn-Regular,
  BoldFont=Vazirmatn-Bold,
  AutoFakeSlant=0.15
]{Vazirmatn}
\setlatintextfont{Latin Modern Roman}

\pagestyle{empty}
\setlength{\parindent}{0pt}
\linespread{1.15}
\hfuzz=2pt

% ---------- کادر مشکی گرد دور تا دور صفحات ----------
\AddToShipoutPictureBG{%
  \begin{tikzpicture}[remember picture,overlay]
    \draw[color=black, line width=0.8pt, rounded corners=3mm] 
      ([xshift=8mm,yshift=-8mm]current page.north west) 
      rectangle 
      ([xshift=-8mm,yshift=8mm]current page.south east);
  \end{tikzpicture}%
}

% ---------- عنوان بخش ----------
\newcommand{\cvsection}[2]{%
  \vspace{16pt}%
  \noindent\LR{\color{primary}\large\faIcon{#1}}\hspace{6pt}%
  {\large\bfseries\color{black}#2}\par
  \vspace{1pt}%
  \noindent\begin{tikzpicture}
    \draw[dashed, lightgray, line width=0.5pt] (0,0) -- (\linewidth,0);
  \end{tikzpicture}\par
  \vspace{8pt}%
}

% ---------- نوار مهارت / زبان ----------
\newcommand{\langbar}[1]{%
  \LR{%
    \begin{tikzpicture}[baseline=-0.6ex]
      \foreach \i in {1,...,5} {
        \ifnum\i>#1\relax
          \fill[lightgray] (\i*0.26,0) rectangle ++(0.20,0.13);
        \else
          \fill[primary]   (\i*0.26,0) rectangle ++(0.20,0.13);
        \fi
      }
    \end{tikzpicture}%
  }%
}

% ---------- بولت‌لیست با tab (۲em) ----------
\newenvironment{bullets}{%
  \begin{itemize}[nosep, leftmargin=2em, itemindent=0pt, itemsep=3pt]%
}{\end{itemize}}

% ---------- توضیح زیرمجموعه‌ی پروژه ----------
\newenvironment{projdesc}{%
  \begin{itemize}[label={}, leftmargin=2em, itemindent=0pt, nosep,
                  itemsep=0pt, topsep=3pt]%
}{\end{itemize}}

% ---------- بلاک اطلاعات با تورفتگی مشابه بولت‌ها ----------
\newenvironment{cvblock}{%
  \begin{itemize}[label={}, leftmargin=2em, itemindent=0pt, nosep, itemsep=6pt, topsep=0pt]%
}{\end{itemize}}

\begin{document}

% ========== هدر ==========
\begin{tcolorbox}[colback=primary, colframe=primary, rounded corners, arc=5mm,
                  boxrule=0pt, left=10pt, right=10pt, top=10pt, bottom=10pt]
  \begin{minipage}[c]{0.72\linewidth}
    \color{headlight}
    {\LARGE\bfseries\color{white}سیدعلی ماجدی}\par
    \vspace{12pt}
    {\small\color{headlight}کارشناس مهندسی صنایع}\par
    \vspace{10pt}
    \small
    \LR{\faTelegram}\ \LR{AliMajedi\_83}\par\vspace{1pt}
    \LR{\faEnvelope}\ \LR{s110majedi@gmail.com}\par\vspace{1pt}
    \LR{\faMapMarker*}\ تهران، خیابان ۱۷ شهریور\par
    \vspace{4pt}
    {\footnotesize متولد: \LR{1383/01/30} \hfill وضعیت تأهل: مجرد \hfill سربازی: معافیت تحصیلی}
  \end{minipage}%
  \hfill
  \begin{minipage}[c]{0.24\linewidth}
    \centering
    \IfFileExists{images/10.jpg}{%
      \begin{tikzpicture}
        \clip (0,0) circle (1.45cm);
        \node at (0,0) {\includegraphics[width=2.9cm]{images/10.jpg}};
      \end{tikzpicture}%
    }{%
      \begin{tikzpicture}
        \fill[white] (0,0) circle (1.45cm);
        \node[color=primary] at (0,0) {\Huge\faUser};
      \end{tikzpicture}%
    }
  \end{minipage}
\end{tcolorbox}

% ========== خلاصه رزومه ==========
\cvsection{user}{خلاصه رزومه}
\begin{bullets}
  \item دانشجوی مهندسی صنایع شریف مسلط به پایتون، لینوکس و علم داده.
  \item علاقه‌مند به حل مسائل پیچیده کامپیوتری در لینوکس، خودکارسازی فرآیندها.
  \item آشنایی و فعالیت در فضاهای ابری و پلتفرم‌های اشتراک‌گذاری کد مانند \LR{GitHub}، \LR{Kaggle} و \LR{Google Colab}.
  \item سابقه توسعه پروژه‌های کاربردی شخصی و روحیه خودآموزی بالا جهت پیوستن به تیم‌های فنی و مهندسی.
\end{bullets}

% ========== سوابق تحصیلی ==========
\cvsection{graduation-cap}{سوابق تحصیلی}
\begin{cvblock}
\item \textbf{\color{primary}کارشناسی مهندسی صنایع}\par
{\small\color{muted}موسسه/دانشگاه: دانشگاه صنعتی شریف --- تهران}
\end{cvblock}

% ========== مهارت‌های نرم ==========
\cvsection{handshake}{مهارت‌های نرم}
\begin{bullets}
  \item مسئولیت‌پذیری، علاقه‌مند به یادگیری مستمر و توانایی حل مسئله.
\end{bullets}

% ========== مهارت‌ها ==========
\cvsection{lightbulb}{مهارت‌ها}
\vspace{4pt}
\noindent
\begin{LTR}
\begin{tabularx}{\dimexpr\linewidth-2em\relax}{X c}
  \RL{\textbf{\color{primary}زبان \LR{MQL5} و الگوریتم‌های معاملاتی (تریدینگ)}} & \langbar{3} \\[4pt]
  \RL{\textbf{\color{primary}مستندسازی با \LR{LaTeX}}} & \langbar{4} \\[4pt]
  \RL{\textbf{\color{primary}ماشین لرنینگ و تحلیل داده}} & \langbar{3} \\[4pt]
  \RL{\textbf{\color{primary}پایگاه داده}} & \langbar{2} \\[4pt]
  \RL{\textbf{\color{primary}پایتون}} & \langbar{4} \\[4pt]
  \RL{\textbf{\color{primary}لینوکس و مفاهیم سیستم عامل}} & \langbar{4} \\[4pt]
  \RL{\textbf{\color{primary}گیت و گیت‌هاب}} & \langbar{3}
\end{tabularx}
\end{LTR}

% ========== زبان ==========
\cvsection{globe}{زبان}
\vspace{4pt}
\noindent
\begin{LTR}
\begin{tabularx}{\dimexpr\linewidth-2em\relax}{X c}
  \RL{\textbf{\color{primary}انگلیسی}} & \langbar{4}
\end{tabularx}
\end{LTR}

% ========== پروژه‌ها ==========
\cvsection{project-diagram}{پروژه‌ها}

% --- پروژه ۱ ---
\begin{cvblock}
\item \textbf{\color{primary}سامانه رزرو بلیت سینما (\LR{CinemaSharif})}\hfill{\small\color{muted}\LR{آذر 1404}}\par
{\small\color{muted}کارفرما / درخواست‌کننده: پروژه شخصی / دانشگاهی}\par
\begin{projdesc}
  \item {\small توسعه سامانه کامل رزرو بلیت و مدیریت سینما با استفاده از \LR{Django} و \LR{PostgreSQL}. پیاده‌سازی سیستم احراز هویت با شماره موبایل، سیستم کیف پول داخلی با گزارش تراکنش‌ها، و پنل مدیریت اختصاصی.}
\end{projdesc}
\vspace{2pt}
\hfill\LR{\href{https://github.com/AliMajedi-83/CinemaSharif}{\textcolor{linkblue}{github.com/AliMajedi-83/CinemaSharif}}}
\end{cvblock}

% --- پروژه ۲ ---
\begin{cvblock}
\item \textbf{\color{primary}ربات اتوماسیون شرکت در کلاس‌های آنلاین (\LR{Skyroom-Auto-Attend})}\par
{\small\color{muted}کارفرما / درخواست‌کننده: پروژه شخصی}\par
\begin{projdesc}
  \item {\small طراحی و توسعه اپلیکیشن پایتون با معماری دوگانه \LR{GUI/CLI} برای ورود خودکار به کلاس‌های آنلاین اسکای‌روم. ضبط خودکار جلسات با استفاده از کتابخانه‌های صوتی پایتون. پیاده‌سازی قابلیت‌هایی نظیر مدیریت جدول زمانی هفتگی، الگوریتم هوشمند و ترکیب خودکار فایل‌های چندرسانه‌ای در پس‌زمینه.}
\end{projdesc}
\vspace{2pt}
\hfill\LR{\href{https://github.com/AliMajedi-83/Skyroom-Auto-Attend}{\textcolor{linkblue}{github.com/AliMajedi-83/Skyroom-Auto-Attend}}}
\end{cvblock}

% --- پروژه ۳ ---
\begin{cvblock}
\item \textbf{\color{primary}استخراج‌کننده کندل‌های بازار مالی (\LR{MT5-M1-BidAsk-Exporter})}\par
{\small\color{muted}کارفرما / درخواست‌کننده: پروژه شخصی}\par
\begin{projdesc}
  \item {\small توسعه ابزار بهینه \LR{Expert Advisor} به زبان \LR{MQL5} جهت اجرا در \LR{Strategy Tester}. قابلیت ثبت و استخراج دقیق آخرین قیمت‌های \LR{Bid} و \LR{Ask} همراه با اطلاعات کامل \LR{OHLC} و \LR{Spread} و کندل‌های ۱ دقیقه‌ای (\LR{M1}) جهت تحلیل در پایتون.}
\end{projdesc}
\vspace{2pt}
\hfill\LR{\href{https://github.com/AliMajedi-83/MT5-M1-BidAsk-Exporter}{\textcolor{linkblue}{github.com/AliMajedi-83/MT5-M1-BidAsk-Exporter}}}
\end{cvblock}

% --- پروژه ۴ ---
\begin{cvblock}
\item \textbf{\color{primary}داشبوردهای تحلیل آماری و کنترل کیفیت مهندسی (\LR{Statistics\_Project})}\par
{\small\color{muted}کارفرما / درخواست‌کننده: پروژه دانشگاهی}\par
\begin{projdesc}
  \item {\small توسعه مجموعه‌ای از داشبوردهای وب و تحلیلی با استفاده از \LR{R} و \LR{Shiny} جهت پیاده‌سازی مفاهیم آمار مهندسی و کنترل کیفیت آماری. تحلیل داده‌های واقعی ساختاری و سنسوری (مانند دیتاست \LR{SECOM})، رسم منحنی‌های مشخصه عملیاتی (\LR{OC Curves}) در طرح‌های نمونه‌برداری، تحلیل‌های \LR{ANOVA}، آزمون‌های فرض و شبیه‌سازی توزیع‌های آماری همراه با رابط کاربری واکنش‌گرا با \LR{bslib} و \LR{ggplot}.}
\end{projdesc}
\vspace{2pt}
\hfill\LR{\href{https://github.com/AliMajedi-83/Statistics_Project}{\textcolor{linkblue}{github.com/AliMajedi-83/Statistics\_Project}}}
\end{cvblock}

% --- پروژه ۵ (ترنسفورمر فارکس) ---
\begin{cvblock}
\item \textbf{\color{primary}پیش‌بینی سطوح حد سود و ضرر در فارکس با ترنسفورمر (\LR{Forex-SL-TP-Prediction})}\par
{\small\color{muted}کارفرما / درخواست‌کننده: پروژه شخصی / پژوهشی}\par
\begin{projdesc}
  \item {\small توسعه مدل یادگیری عمیق مبتنی بر معماری ترنسفورمر (\LR{Transformer}) جهت پیش‌بینی ماتریس هم‌زمان برخورد قیمت به سطوح حد سود و حد ضرر در بازار فارکس. پیاده‌سازی بلوک‌های \LR{Cross-Attention} برای خلاصه‌سازی ویژگی‌های زمانی کندل‌ها و طراحی تابع خطای اختصاصی \LR{BCE} وزن‌دار و امتیازدهی ریسک‌آگاه با آستانه دینامیک.}
\end{projdesc}
\vspace{2pt}
\hfill\LR{\href{https://colab.research.google.com/drive/11UVGcd8ofB5Pry0TlUdaY8cIlkfylQVu?usp=sharing}{\textcolor{linkblue}{Google Colab}}}
\end{cvblock}

% --- پروژه ۶ (دیتاست کگل) ---
\begin{cvblock}
\item \textbf{\color{primary}دیتاست تقویم اقتصادی فارکس فکتوری (\LR{ForexFactory-News-Dataset})}\par
{\small\color{muted}کارفرما / درخواست‌کننده: پروژه شخصی}\par
\begin{projdesc}
  \item {\small توسعه ربات استخراج وب (\LR{Web Scraper}) جهت جمع‌آوری، پاک‌سازی و ساختاردهی دیتاست جامع اخبار و رویدادهای تقویم اقتصادی \LR{Forex Factory} از سال ۲۰۰۷ تا ۲۰۲۶. انتشار این دیتاست در پلتفرم کگل (\LR{Kaggle}) جهت استفاده پژوهشگران در مدل‌های معاملاتی مبتنی بر تحلیل بنیادی.}
\end{projdesc}
\vspace{2pt}
\hfill\LR{\href{https://www.kaggle.com/datasets/seyyedalimajedi/forexfactory-news-csv}{\textcolor{linkblue}{kaggle.com/datasets/seyyedalimajedi/forexfactory-news-csv}}}
\end{cvblock}

% ========== سوابق شغلی ==========
\cvsection{briefcase}{سوابق شغلی}
\begin{cvblock}
\item \textbf{\color{primary}آموزش ریاضی دبیرستان}\par
{\small\color{muted}مدرسه علوی --- تهران}\par
\vspace{2pt}
{\small\color{muted}وظایف و دستاوردها:}
\begin{bullets}
  \item دستیار معلم برای رفع اشکال و تست‌زنی سال کنکور
\end{bullets}
\end{cvblock}

\end{document}
